\chapter{Preliminaries}

\section{Continuous-Time Stochastic Processes}

We introduce several basic notions for continuous-time stochastic processes.

\begin{de}
A stochastic process is a family
\[
X=(X_t)_{t\in\mathbb T}
\]
of random variables taking values in a measurable space \(\mathcal X\) and indexed by time \(t\). The time set \(\mathbb T\) may be discrete or continuous.
\end{de}

\begin{remark}
For the purposes of this book, we take \(\mathcal X=\mathbb R^d\), equipped with its Borel \(\sigma\)-field, and consider continuous-time processes. The time interval is either finite, \(\mathbb T=[0,T]\) with \(0<T<\infty\), or infinite, \(\mathbb T=[0,\infty)\).
\end{remark}

We next introduce filtrations and measurability.

\begin{de}
A filtration on \((\Omega,\mathcal F,P)\) is an increasing family
\[
\mathbb F=(\mathcal F_t)_{t\in\mathbb T}
\]
of sub-\(\sigma\)-fields of \(\mathcal F\); that is,
\[
\mathcal F_s\subseteq\mathcal F_t\subseteq\mathcal F,
\qquad 0\le s\le t,
\quad s,t\in\mathbb T.
\]
\end{de}

\begin{de}
Let
\[
\mathcal N^P:=\{N\subseteq A:A\in\mathcal F,\ P(A)=0\}
\]
be the collection of \(P\)-negligible sets. A filtration \(\mathbb F\) satisfies the usual conditions if it is complete,
\[
\mathcal N^P\subseteq\mathcal F_0,
\]
and right-continuous. For \(t<\sup\mathbb T\), right-continuity means
\[
\mathcal F_{t+}
:=\bigcap_{\substack{s>t\\s\in\mathbb T}}\mathcal F_s
=\mathcal F_t.
\]
On a finite horizon, we use the convention \(\mathcal F_{T+}=\mathcal F_T\).
\end{de}

\begin{remark}
Right-continuity means, informally, that observing the filtration immediately after time \(t\) reveals no additional information beyond that already available at time \(t\). Completeness means that every subset of a null event is already known at time zero.
\end{remark}

\begin{remark}
Starting from an arbitrary filtration, one obtains its usual augmentation by adjoining all \(P\)-negligible sets and then imposing right-continuity. On an infinite horizon one may write
\[
\mathcal F_t^{\mathrm{aug}}
:=\bigcap_{s>t}\bigl(\mathcal F_s\vee\mathcal N^P\bigr),
\]
and on a finite horizon one uses the analogous definition for \(t<T\), together with \(\mathcal F_T^{\mathrm{aug}}=\mathcal F_T\vee\mathcal N^P\).
\end{remark}

A process can be viewed either pathwise, as a function of time for fixed \(\omega\), or as a random variable \(X_t\) for fixed \(t\). Considering the joint map \((t,\omega)\mapsto X_t(\omega)\) leads to the following definitions.

\begin{de}[Progressively measurable, optional, and predictable processes]
A process \(X\) is \emph{progressively measurable} if, for every \(t\in\mathbb T\), its restriction to \([0,t]\times\Omega\) is measurable with respect to the product \(\sigma\)-field \(\mathcal B([0,t])\otimes\mathcal F_t\).

A process is \emph{optional} if \((t,\omega)\mapsto X_t(\omega)\) is measurable with respect to the optional \(\sigma\)-field, namely the \(\sigma\)-field generated by all adapted c\`adl\`ag processes.

A process is \emph{predictable} if \((t,\omega)\mapsto X_t(\omega)\) is measurable with respect to the predictable \(\sigma\)-field, namely the \(\sigma\)-field generated by all adapted left-continuous processes.
\end{de}

We shall often study processes through their conditional-expectation behavior.

\begin{de}
An adapted process \(X=(X_t)_{t\in\mathbb T}\) is a supermartingale if \(X_t\in L^1\) for every \(t\in\mathbb T\) and
\[
\mathbb E[X_t\mid\mathcal F_s]\le X_s
\quad\text{a.s.},
\qquad 0\le s\le t,
\quad s,t\in\mathbb T.
\]
It is a submartingale if \(-X\) is a supermartingale, and it is a martingale if it is both a supermartingale and a submartingale.
\end{de}

The martingale notion can be localized.

\begin{de}
Let \(X\) be an adapted c\`adl\`ag process. It is a local martingale if there is an increasing sequence of stopping times \((\tau_n)_{n\ge1}\) such that \(\tau_n\uparrow\infty\) a.s.\ and each stopped process \(X^{\tau_n}\) is a martingale. On a finite horizon \([0,T]\), one equivalently requires \(\tau_n\uparrow T\) a.s.
\end{de}

Local martingales form the natural setting for stochastic integration. The following criterion is frequently useful for proving that a local martingale is a true martingale.

\begin{prop}
Let \(M=(M_t)_{t\in\mathbb T}\) be a local martingale. If
\[
\mathbb E\!\left[\sup_{0\le s\le t}|M_s|\right]<\infty
\qquad\text{for every }t\in\mathbb T,
\]
then \(M\) is a martingale.
\end{prop}

\begin{remark}
A c\`adl\`ag local martingale is a martingale if and only if it is locally of class \((DL)\): for every deterministic \(a<\infty\), the family
\[
\{M_\tau:\tau\le a,\ \tau\text{ a stopping time}\}
\]
is uniformly integrable. On a fixed finite horizon this is the usual class-\((D)\) condition. The preceding proposition is a convenient sufficient criterion.
\end{remark}

We now summarize the Snell envelope, which is fundamental in optimal stopping. For \(0\le t\le T<\infty\), let \(\mathcal T_{t,T}\) denote the set of stopping times taking values in \([t,T]\).

\begin{prop}[Snell envelope]
Let \(H=(H_t)_{0\le t\le T}\) be a real-valued, adapted, c\`adl\`ag process of class \((D)\). Its Snell envelope is
\[
V_t
:=\operatorname*{ess\,sup}_{\tau\in\mathcal T_{t,T}}
\mathbb E[H_\tau\mid\mathcal F_t],
\qquad 0\le t\le T.
\]
Then \(V\) is the smallest c\`adl\`ag supermartingale dominating \(H\).

If, in addition, the filtration is continuous and \(H\) has no negative jumps,
\[
H_t-H_{t-}\ge0,
\qquad 0<t\le T,
\]
then \(V\) is continuous and
\[
\tau_t^*
:=\inf\{s\in[t,T]:V_s=H_s\}\wedge T
\]
is optimal after time \(t\). In particular,
\[
V_t
=\mathbb E[V_{\tau_t^*}\mid\mathcal F_t]
=\mathbb E[H_{\tau_t^*}\mid\mathcal F_t].
\]
\end{prop}

\section{Stochastic Calculus}

We record several standard tools from stochastic calculus.

\begin{thm}[It\^o's formula]
Let \(X\) be an \(n\)-dimensional It\^o process satisfying
\[
\mathrm dX_t=\mu_t\,\mathrm dt+\sigma_t\,\mathrm dW_t,
\]
where \(\mu_t\in\mathbb R^n\), \(\sigma_t\in\mathbb R^{n\times m}\), and \(W\) is an \(m\)-dimensional Brownian motion. If \(f\in C^{1,2}([0,T]\times\mathbb R^n)\), then
\[
\begin{aligned}
\mathrm df(t,X_t)
={}&\left[
\partial_t f(t,X_t)
+\nabla_x f(t,X_t)^\top\mu_t
+\frac12\operatorname{tr}\!\left(
\sigma_t\sigma_t^\top D_x^2f(t,X_t)
\right)
\right]\mathrm dt\\
&+\nabla_x f(t,X_t)^\top\sigma_t\,\mathrm dW_t.
\end{aligned}
\]
\end{thm}

We next introduce the infinitesimal generator associated with a diffusion.

\begin{de}
Suppose that \(X\) solves
\[
\mathrm dX_t=\mu(t,X_t)\,\mathrm dt+\sigma(t,X_t)\,\mathrm dW_t.
\]
For a twice continuously differentiable test function \(\varphi\), the time-dependent generator is
\[
(\mathcal A_t\varphi)(x)
:=\sum_{i=1}^n\mu_i(t,x)\,\partial_{x_i}\varphi(x)
+\frac12\sum_{i,j=1}^n
(\sigma\sigma^\top)_{ij}(t,x)\,
\partial_{x_ix_j}^2\varphi(x).
\]
For a space-time test function \(F\), the corresponding space-time operator is \(\partial_tF+\mathcal A_tF\).
\end{de}

The Feynman--Kac formula links diffusion processes with parabolic partial differential equations.

\begin{thm}[Feynman--Kac formula]
Let \(X^{t,x}\) be the solution of
\[
\begin{aligned}
\mathrm dX_s^{t,x}
&=\mu(s,X_s^{t,x})\,\mathrm ds
+\sigma(s,X_s^{t,x})\,\mathrm dW_s,
\qquad t\le s\le T,\\
X_t^{t,x}&=x.
\end{aligned}
\]
Assume that \(F\in C^{1,2}([0,T]\times\mathbb R^n)\) has sufficient growth control for It\^o's formula and solves
\[
\begin{aligned}
\partial_tF(t,x)+(\mathcal A_tF)(t,x)-rF(t,x)&=0,
&&0\le t<T,\\
F(T,x)&=\Phi(x).
\end{aligned}
\]
Assume also that
\[
\mathbb E_{t,x}\!\left[
\int_t^T e^{-2r(s-t)}
\left\|
\sigma(s,X_s^{t,x})^\top\nabla_xF(s,X_s^{t,x})
\right\|^2\mathrm ds
\right]<\infty.
\]
Then
\[
F(t,x)
=e^{-r(T-t)}\mathbb E_{t,x}[\Phi(X_T^{t,x})],
\]
where \(\mathbb E_{t,x}\) denotes expectation conditional on \(X_t=x\).
\end{thm}

Transition densities satisfy the Kolmogorov backward and forward equations.

\begin{thm}[Kolmogorov backward equation]
Suppose that the diffusion above has a transition density \(p\) defined by
\[
P(s,y;t,\mathrm dx)
:=P(X_t\in\mathrm dx\mid X_s=y)
=p(s,y;t,x)\,\mathrm dx.
\]
Under the regularity assumptions needed to differentiate the density,
\[
\left(\partial_s+\mathcal A_s^{(y)}\right)p(s,y;t,x)=0,
\qquad 0\le s<t,
\]
where \(\mathcal A_s^{(y)}\) acts on the starting variable \(y\). The terminal condition is
\[
p(s,\cdot\,;t,x)\longrightarrow\delta_x
\quad\text{as }s\uparrow t
\]
in the sense of distributions in the variable \(y\).
\end{thm}

\begin{thm}[Kolmogorov forward equation]
Under the same regularity assumptions, the density also satisfies
\[
\partial_t p(s,y;t,x)
=\mathcal A_t^*p(s,y;t,x),
\qquad t>s,
\]
with initial condition
\[
p(s,y;s,\cdot)\,\mathrm dx=\delta_y.
\]
Writing
\[
a(t,x):=\sigma(t,x)\sigma(t,x)^\top,
\]
the formal adjoint is
\[
(\mathcal A_t^*q)(x)
=-\sum_{i=1}^n\partial_{x_i}\bigl(\mu_i(t,x)q(x)\bigr)
+\frac12\sum_{i,j=1}^n
\partial_{x_ix_j}^2\bigl(a_{ij}(t,x)q(x)\bigr).
\]
This is also called the Fokker--Planck equation.
\end{thm}

\section{Regularity Estimates in Stochastic Analysis}

We begin with Doob's maximal inequalities.

\begin{thm}[Doob's maximal inequalities]
Let \(\tau\) be a bounded stopping time.
\begin{enumerate}[label=\textup{(\roman*)}]
\item If \(X\) is a nonnegative c\`adl\`ag submartingale, then for every \(\lambda>0\),
\[
P\!\left(\sup_{0\le s\le\tau}X_s\ge\lambda\right)
\le\frac{\mathbb E[X_\tau]}{\lambda}.
\]

\item If \(X\) is a c\`adl\`ag martingale, \(p>1\), and \(X_\tau\in L^p\), then
\[
\mathbb E\!\left[
\left(\sup_{0\le s\le\tau}|X_s|\right)^p
\right]
\le
\left(\frac{p}{p-1}\right)^p
\mathbb E[|X_\tau|^p].
\]
\end{enumerate}
\end{thm}

\begin{remark}
If \(X\) is a c\`adl\`ag uniformly integrable martingale on an infinite horizon, then the martingale convergence theorem and Doob's inequality imply
\[
\sup_{t\ge0}|X_t|<\infty
\quad\text{a.s.}
\]
\end{remark}

The following inequality controls covariation and is used in stochastic integration.

\begin{prop}[Kunita--Watanabe inequality]
Let \(M\) and \(N\) be continuous local martingales, and let \(\alpha\) and \(\beta\) be predictable processes for which the integrals below are well defined. Then, for every \(t\),
\[
\int_0^t|\alpha_s\beta_s|\,
\mathrm d|\langle M,N\rangle|_s
\le
\left(\int_0^t\alpha_s^2\,\mathrm d\langle M\rangle_s\right)^{1/2}
\left(\int_0^t\beta_s^2\,\mathrm d\langle N\rangle_s\right)^{1/2}
\quad\text{a.s.}
\]
\end{prop}

\begin{remark}
Set \(A:=\langle M\rangle+\langle N\rangle\). There are predictable densities \(a,b,c\) such that
\[
\mathrm d\langle M\rangle=a\,\mathrm dA,
\qquad
\mathrm d\langle N\rangle=b\,\mathrm dA,
\qquad
\mathrm d\langle M,N\rangle=c\,\mathrm dA,
\]
and the Kunita--Watanabe inequality yields \(|c|\le\sqrt{ab}\) almost everywhere with respect to \(\mathrm dA\).
\end{remark}

\begin{thm}[Burkholder--Davis--Gundy inequality]
For every \(p>0\), there are constants \(c_p,C_p>0\) such that, for every continuous local martingale \(M\) with \(M_0=0\) and every stopping time \(\tau\),
\[
c_p\,\mathbb E[\langle M\rangle_\tau^{p/2}]
\le
\mathbb E\!\left[
\left(\sup_{0\le s\le\tau}|M_s|\right)^p
\right]
\le
C_p\,\mathbb E[\langle M\rangle_\tau^{p/2}].
\]
The inequalities are understood in the extended sense when one of the expectations is infinite.
\end{thm}

\begin{cor}
Let \(M\) be a continuous local martingale with \(M_0\in L^1\). If
\[
\mathbb E[\langle M\rangle_t^{1/2}]<\infty
\qquad\text{for every finite }t,
\]
then \(M\) is a martingale.
\end{cor}

We next record standard moment and stability estimates for stochastic differential equations.

\begin{thm}[Existence, uniqueness, and a moment estimate]
Consider the \(\mathbb R^n\)-valued SDE
\[
\mathrm dX_s=b(s,X_s,\omega)\,\mathrm ds
+\sigma(s,X_s,\omega)\,\mathrm dW_s.
\]
Assume that, for each \(x\), the coefficient processes are progressively measurable and that there are a deterministic constant \(K\ge0\) and a nonnegative progressively measurable process \(\kappa\) such that
\[
\begin{aligned}
|b(s,x)-b(s,y)|+\|\sigma(s,x)-\sigma(s,y)\|
&\le K|x-y|,\\
|b(s,x)|+\|\sigma(s,x)\|
&\le\kappa_s+K|x|.
\end{aligned}
\]
Fix \(p\ge2\) and a finite horizon \(T\), and assume
\[
\mathbb E\!\left[
\left(\int_0^T\kappa_s^2\,\mathrm ds\right)^{p/2}
\right]<\infty.
\]
Then, for every \(t\le T\) and every \(\mathcal F_t\)-measurable \(\xi\in L^p(\Omega;\mathbb R^n)\), there is a unique strong solution satisfying \(X_t=\xi\), and
\[
\mathbb E\!\left[
\sup_{t\le s\le T}|X_s|^p
\right]
\le C_{p,T}
\left(
1+\mathbb E|\xi|^p
+\mathbb E\!\left[
\left(\int_t^T\kappa_s^2\,\mathrm ds\right)^{p/2}
\right]
\right).
\]
\end{thm}

\begin{thm}[Stability estimates]
Under the preceding assumptions, let \(X^{t,x}\) denote the solution starting from the deterministic point \(x\) at time \(t\). For every finite \(T\), there is a constant \(C_T\) such that, for \(t\le s\le T\),
\[
\mathbb E\!\left[
\sup_{t\le u\le s}|X_u^{t,x}|^2
\right]
\le
C_T\left(
|x|^2+
\mathbb E\int_t^s\kappa_r^2\,\mathrm dr
\right),
\]
and
\[
\mathbb E\!\left[
\sup_{t\le u\le s}|X_u^{t,x}-x|^2
\right]
\le
C_T\,
\mathbb E\int_t^s\bigl(|x|^2+\kappa_r^2\bigr)\,\mathrm dr.
\]

If, in addition, there is a constant \(\beta_0\in\mathbb R\) such that
\[
\begin{aligned}
&(b(r,x)-b(r,y))\cdot(x-y)\\
&\qquad
+\frac12\|\sigma(r,x)-\sigma(r,y)\|^2
\le\beta_0|x-y|^2
\end{aligned}
\]
for all \(r,x,y\), almost surely, then
\[
\mathbb E|X_s^{t,x}-X_s^{t,y}|^2
\le e^{2\beta_0(s-t)}|x-y|^2.
\]
Moreover, without strengthening the exponential constant to \(e^{2\beta_0(s-t)}\), one has the supremum estimate
\[
\mathbb E\!\left[
\sup_{t\le u\le s}|X_u^{t,x}-X_u^{t,y}|^2
\right]
\le C_T|x-y|^2.
\]
\end{thm}

\begin{proof}
Apply It\^o's formula to \(|X_s|^2\), \(|X_s-x|^2\), or \(|X_s^{t,x}-X_s^{t,y}|^2\), as appropriate. For estimates involving a supremum, use the Burkholder--Davis--Gundy inequality on the stochastic integral and then Young's inequality to absorb the resulting supremum term. The remaining deterministic integral inequality has the form
\[
g(t)\le h(t)+C\int_0^tg(s)\,\mathrm ds,
\]
so Gronwall's lemma yields the first two estimates and the supremum stability bound. Under the one-sided monotonicity condition, taking expectations before introducing a supremum gives
\[
\mathbb E|X_s^{t,x}-X_s^{t,y}|^2
\le |x-y|^2
+2\beta_0\int_t^s
\mathbb E|X_r^{t,x}-X_r^{t,y}|^2\,\mathrm dr,
\]
and a second application of Gronwall's lemma gives the endpoint estimate.
\end{proof}
